\subsection{مبانی نظری و فرمول‌بندی ریاضی}
فرض کنید $\mathbf{X} \in \mathbb{R}^{B \times N \times D}$ ماتریس حالت‌های پنهان ورودی لایه توجه باشد که در آن $B$ اندازه دسته، $N$ طول دنباله و $D$ بعد پنهان مدل است. در لایه‌های نگاشت خطی \lr{FactorizedLinear}، ماتریس وزن کامل $\mathbf{W} \in \mathbb{R}^{D \times D_{\text{out}}}$ از طریق ضرب دو ماتریس رتبه تقلیل‌یافته $r \ll D$ به همراه نگاشت تطبیقی موازی مدل‌سازی می‌شود:
\begin{equation}
\mathbf{Y} = \mathbf{X} \mathbf{W}_1 \mathbf{W}_2 + \frac{\alpha}{r_{\text{lora}}} \mathbf{X} \mathbf{A}_{\text{lora}} \mathbf{B}_{\text{lora}}
\end{equation}
که در آن $\mathbf{W}_1 \in \mathbb{R}^{D \times r}$ و $\mathbf{W}_2 \in \mathbb{R}^{r \times D_{\text{out}}}$ ماتریس‌های فاکتورسازی ساختاری بوده و $\mathbf{A}_{\text{lora}} \in \mathbb{R}^{D \times r_{\text{lora}}}$ و $\mathbf{B}_{\text{lora}} \in \mathbb{R}^{r_{\text{lora}} \times D_{\text{out}}}$ پارامترهای آموزش‌پذیر باقیمانده هستند.

تانسورهای پرس‌وجو، کلید و مقدار پس از تفکیک به $H$ سر توجه و بعد هر سر $d = D/H$ به دست می‌آیند:
\begin{equation}
\mathbf{Q} \in \mathbb{R}^{B \times H \times N \times d}, \quad \mathbf{K}, \mathbf{V} \in \mathbb{R}^{B \times H_{\text{kv}} \times N \times d}
\end{equation}
خروجی لایه از تلفیق محدب شاخه محلی و سراسری تشکیل می‌گردد:
\begin{equation}
\mathbf{O}_{\text{local}} = \operatorname{LocalWindowAttn}(\mathbf{Q}, \mathbf{K}, \mathbf{V}, W)
\end{equation}
در شاخه سراسری، ابتدا توابع نگاشت کِرنل غیرخطی $\phi(\cdot) = \operatorname{ELU}(\cdot) + 1$ اعمال شده و ضریب هرس تطبیقی $\mathbf{g}_{\text{prune}} = \sigma(\mathbf{X} \mathbf{W}_{\text{prune}})$ برای کنترل چگالی اطلاعاتی در کلیدها ضرب می‌شود:
\begin{equation}
\tilde{\mathbf{K}} = \phi(\mathbf{K}) \odot \mathbf{g}_{\text{prune}}
\end{equation}
در حالت علّی، توجه خطی با جمع تجمعی تکه‌ای بر روی بعد زمان محاسبه می‌گردد:
\begin{equation}
\mathbf{S}_t = \sum_{\tau=1}^t \tilde{\mathbf{K}}_\tau^\top \mathbf{V}_\tau, \quad \mathbf{Z}_t = \sum_{\tau=1}^t \tilde{\mathbf{K}}_\tau^\top, \quad \mathbf{O}_{\text{global}, t} = \frac{\phi(\mathbf{Q}_t) \mathbf{S}_t}{\phi(\mathbf{Q}_t) \mathbf{Z}_t + \epsilon}
\end{equation}
در نهایت، گیت همجوشی یادگیرنده $\boldsymbol{\alpha} = \sigma(\mathbf{X} \mathbf{W}_{\text{fuse}}) \in [0, 1]$ تعادل میان توجه محلی و سراسری را برقرار می‌سازد:
\begin{equation}
\mathbf{O}_{\text{fused}} = (1 - \boldsymbol{\alpha}) \odot \mathbf{O}_{\text{local}} + \boldsymbol{\alpha} \odot \mathbf{O}_{\text{global}}
\end{equation}

که در آن $\sigma_i(\mathbf{W})$ بیانگر $i$-امین مقدار منفرد مرتب‌شده ماتریس به صورت نزولی است. نسبت حفظ انرژی طیفی $\eta_r = \frac{\sum_{i=1}^r \sigma_i^2}{\sum_{i=1}^{d} \sigma_i^2}$ تضمین می‌کند که با انتخاب آستانه $\tau \ge 0.90$، بیش از نود درصد از پراکندگی واریانس داده‌ها حفظ شده و افت سیگنال در انتشار گرادیان رو به عقب به حداقل ممکن کاهش می‌یابد.

\paragraph{تحلیل مجانبی و پیچیدگی اشتراک حافظه سخت‌افزاری}
از منظر پیچیدگی محاسباتی و انتقال داده‌ها در سلسومراتب حافظه (حافظه سراسری \lr{HBM} به حافظه محلی \lr{SRAM})، عملیات ضرب ماتریسی در این مکانیزم نیازمند تخصیص بافر موقت با ابعاد فاکتورگیری‌شده است. برای پردازش دنباله‌ای به طول $N$ با بعد مخفی $D$، بار محاسباتی هسته از فرم درجه دو به رابطه زیر تقلیل می‌یابد:
\begin{equation}
\label{eq:complexity_aegis_hybrid}
\mathcal{C}_{\text{aegis\_hybrid}} = \mathcal{O}\left( B \cdot N \cdot D \cdot r + B \cdot H \cdot N \cdot d_k \right) \ll \mathcal{O}\left( B \cdot H \cdot N^2 \cdot d_k \right)
\end{equation}
این بهبود مرتبه پیچیدگی، با کاهش چشمگیر حجم تبادل تانسورها در گذرگاه‌های ارتباطی کارت گرافیک، مانع از اشباع پهنای باند حافظه (پدیده \lr{Memory-Bound}) در فرآیند استنتاج طولانی‌مدت می‌گردد.
